\documentclass[10pt,a4paper]{amsart}
\usepackage[margin=19mm]{geometry}
\usepackage{amsmath,amssymb,mathtools}
\usepackage[scr=rsfs,cal=euler]{mathalfa}
\usepackage{xfrac}
\usepackage[unicode,hidelinks,bookmarks=false]{hyperref}
\newcommand{\ie}{i.e.\ }
\newcommand{\cf}{cf.\ }
\newcommand{\resp}{respectively\ }
\newcommand{\Subsection}{\subsection}
\providecommand{\nobreakdash}{\nobreak\hbox{-}}
\setlength{\emergencystretch}{2em}
\multlinegap0pt
\usepackage{bm}
\usepackage{bbm}
\usepackage{dsfont}
\usepackage{xspace}
\usepackage{cancel}
\usepackage{mathtools}
\usepackage{amsthm}
\makeatletter
\@ifundefined{theo}{\newtheorem{theo}{Theo}}{}
\@ifundefined{defi}{\newtheorem{defi}[theo]{Definition}}{}
\@ifundefined{lemm}{\newtheorem{lemm}[theo]{Lemma}}{}
\makeatother
\title[Existence and uniqueness of the global conservative solution]{Existence and uniqueness of the global conservative solutions for the generalized Camassa-Holm equation with dual-power nonlinearities}
\author{Jian Chen, Xiaoxin Chen, Zhaoyang Yin}
\date{Source: arXiv:2603.12978v2 (CC BY 4.0)}
\begin{document}
\maketitle

\noindent\emph{Source and licence.} Existence and uniqueness of the global conservative solutions for the generalized Camassa-Holm equation with dual-power nonlinearities. Version arXiv:2603.12978v2 (CC BY 4.0), distributed under the Creative Commons Attribution 4.0 International licence, as recorded in the arXiv licence metadata for this exact version and shown on the abstract page https://arxiv.org/abs/2603.12978v2. Full rights notice: licenses/arxiv-2603.12978-CC-BY-4.0.txt.\par\smallskip

\noindent\textit{Editorial scope.} Counted unit: the Theo whose source label is \texttt{None} in the article (source0.tex lines 503--505, source label \texttt{None}), delivered together with the article's own complete original proof of it (source0.tex lines 506--575, one \texttt{proof} environment of 70 lines). No mathematics is written, repaired or re-derived: statement and proof are the article's own text line for line. Required context retained uncounted: eq1 (lines 115--122); eqw (lines 151--153); solution (lines 156--164); u (lines 160--162); conservative (lines 166--175); y0 (lines 187--189); Px (lines 208--213); semi-linear system (lines 230--237); semi-linear system\_initial data (lines 239--246); u\_xi (lines 351--353); bdd\_u (lines 381--383); partial\_t y(t,xi) (lines 401--403); y\_xi (lines 413--415); bdd\_Px\_Linfty (lines 433--435); bdd\_y(t,xi) (lines 449--451); lemma\_N (lines 488--493). Every definition, display and label of those lines is retained unchanged. Omitted from this package and not redistributed: 1--114, 123--150, 154--155, 163--165, 176--186, 190--207, 214--229, 238--238, 247--350, 354--380, 384--400, 404--412, 416--432, 436--448, 452--487, 494--502, 576--1002. Nothing inside a retained excerpt is omitted or altered: every retained body is delivered exactly as the article's lines are written, so no omission receipt of any kind accompanies this package. Citations: the retained excerpts carry no citation command at all, so no bibliography entry of the article is transcribed and no reference is redistributed. Reference adaptation: none was needed and none was made; every reference inside the delivered unit resolves inside it, so no unresolved reference or citation is produced. Printed slips kept literal: no printed word, symbol, hypothesis or number of the counted statement or of its proof was altered. Layout and class: the article's own class is \texttt{article} with packages including amsmath, amssymb, amsthm, mathrsfs, dsfont, cite, geometry, titlesec, hyperref, tocloft, color, fancyhdr; the wrapper is the lane's pinned \texttt{amsart} editorial wrapper at 10pt on A4, which restates the theorem-like environments the retained text needs and adds the article's own macro definitions that the retained text needs (none), with extra wrapper packages (bm, bbm, dsfont, xspace, cancel, mathtools). The delivered numbering is the unit's own. Assets: no retained span contains a figure environment, a graphics-inclusion command, a PDF-page inclusion, a table or a TikZ picture; no image file is copied into this package, the package declares no asset dependency and no image file is redistributed with it. Licence: version arXiv:2603.12978v2 of this article is distributed under the Creative Commons Attribution 4.0 International licence, as recorded in the arXiv licence metadata for this exact version and shown on the abstract page https://arxiv.org/abs/2603.12978v2. A full rights notice is delivered at licenses/arxiv-2603.12978-CC-BY-4.0.txt. Characters: every retained excerpt is pure ASCII, so the delivered engine, XeLaTeX with its default Latin Modern text font, reports no missing character.
	\begin{equation}\label{eq1}
		\left\{
		\begin{array}{ll}
			u_t+\left(\frac{s}{\mu}u-\frac{\eta}{\mu}\right)u_x
			=-P_x,\ &t>0,\ x\in\mathbb{R}, \\
			u(0,x)=u_0(x), &x\in\mathbb{R},
		\end{array}\right.
	\end{equation}

	\begin{equation}\label{eqw}
		w_t+\left(\left(\frac{s}{\mu}u-\frac{\eta}{\mu}\right)w\right)_x=\frac{2}{\mu}u_x\left(-\left(\frac A2+\frac{s}{2\mu}\right)u^2-\frac{B}{m+1}u^{m+1}+\left(2k+\frac{\eta}{\mu}\right)u-P\right).
	\end{equation}

	\begin{defi}\label{solution}
		We call $u=u(t,x)$ a solution of the Cauchy problem \eqref{eq1} on $[0,T]$ if $u$ is a H\"{o}lder continuous function defined on $[0,T]\times\mathbb{R}$ with the following properties. \\
		1. $u(t,\cdot)\in H^1(\mathbb{R}),\ \forall\ t\in[0,T]$.\\
		2. the map $t\mapsto u(t,\cdot)$ is Lipschitz continuous from $[0,T]$ into $L^2(\mathbb{R})$ and satisfies the initial condition $u(0,x)=u_0(x)$ together with
		\begin{equation}\label{u}
			\frac{d}{dt}u=-\left(\frac{s}{\mu}u-\frac{\eta}{\mu}\right)u_x-P_x
		\end{equation}
		for a.e. $t$. Here \eqref{u} is understood as an equality between functions in $L^2(\mathbb{R})$.
	\end{defi}

	\begin{defi}\label{conservative}
		We call $u=u(t,x)$ a global conservative solution of \eqref{eq1} if $u$ satisfies the following properties. \\
		1. $u$ provides a solution to the Cauchy problem \eqref{eq1} on $[0,\infty)$ in sense of Definition \ref{solution}.\\
		2. There exists a family of bounded Radon measures $\left\{\mu_{(t)},\ t\geq0\right\}$, depending continuously on time with respect to the topology of weak convergence of measures, whose absolutely continuous part has density $u_x^2$ with respect to the Lebesgue measure. This family provides a measure-valued solution $w$ to the balance law \eqref{eqw}, namely
		\begin{align}
			&\int_{0}^{\infty}\int_\mathbb{R}\varphi_t+\varphi_x\left(\frac{s}{\mu}u-\frac{\eta}{\mu}\right)d\mu_{(t)}dt+\int_{\mathbb{R}}u_{0x}^{2}\varphi(0,x)dx\notag\\
			+&\frac{2}{\mu}\int_{0}^{\infty}\int_\mathbb{R}\left(-\left(\frac{A}{2}+\frac{s}{2\mu}\right)u^2-\frac{B}{m+1}u^{m+1}+\left(2k+\frac{\eta}{\mu}\right)u-P\right)u_x\varphi dxdt=0\label{w}
		\end{align}
		for every test function $\varphi\in\mathcal{C}_c^1(\mathbb{R}^+\times\mathbb{R})$.
	\end{defi}

	\begin{equation}\label{y0}
		\int_0^{y_0(\xi)}(1+u_{0x}^2)dx=\xi.
	\end{equation}

	\begin{align}
		P(\xi) =&\frac{1}{2\sqrt{\mu}}\int_{-\infty}^{\infty}\exp\left\{-\frac{1}{\sqrt{\mu}}\left|\int_{\xi}^{\xi^{\prime}}\cos^{2}\frac{v(s)}{2}\cdot q(s)ds\right|\right\}\notag \\
		& \cdot\left\{\left[-(\frac{A}{2}+\frac{s}{2\mu})u^2(\xi^{\prime})-\frac{B}{m+1}u^{m+1}(\xi^{\prime})+(2k+\frac{\eta}{\mu})u(\xi^{\prime})\right]\cos^2\frac{v(\xi^{\prime})}{2}+\frac{s}{2}\sin^2\frac{v(\xi^{\prime})}{2}\right\}q(\xi^{\prime})d\xi^{\prime},\label{P}\\
		P_x(\xi) =&\frac{1}{2\mu}\left(\int_\xi^\infty-\int_{-\infty}^\xi\right)\exp\left\{-\frac{1}{\sqrt{\mu}}\left|\int_{\xi}^{\xi^{\prime}}\cos^{2}\frac{v(s)}{2}\cdot q(s)ds\right|\right\}\notag \\
		& \cdot\left\{\left[-(\frac{A}{2}+\frac{s}{2\mu})u^2(\xi^{\prime})-\frac{B}{m+1}u^{m+1}(\xi^{\prime})+(2k+\frac{\eta}{\mu})u(\xi^{\prime})\right]\cos^2\frac{v(\xi^{\prime})}{2}+\frac{s}{2}\sin^2\frac{v(\xi^{\prime})}{2}\right\}q(\xi^{\prime})d\xi^{\prime}.\label{Px}
	\end{align}

	\begin{equation}\label{semi-linear system}
		\left\{
		\begin{array}{l}
			\frac{\partial u}{\partial t}=-P_x,\\
			\frac{\partial v}{\partial t}=\frac{1}{\mu}(1+\cos v)\left(-\left(\frac A2+\frac{s}{2\mu}\right)u^2-\frac{B}{m+1}u^{m+1}+\left(2k+\frac{\eta}{\mu}\right)u-P\right)-\frac{s}{\mu}\sin^2\frac{v}{2},\\
			\frac{\partial q}{\partial t}=\frac{1}{\mu}\left(\frac{s}{2}-\left(\frac A2+\frac{s}{2\mu}\right)u^2-\frac{B}{m+1}u^{m+1}+\left(2k+\frac{\eta}{\mu}\right)u-P\right)\sin v\cdot q,
		\end{array}\right.
	\end{equation}

	\begin{equation}\label{semi-linear system_initial data}
		\left\{
		\begin{array}{l}
			u(0,\xi)=u_0\left(y_0(\xi)\right), \\
			v(0,\xi)=2\arctan u_{0x}\left(y_0(\xi)\right), \\
			q(0,\xi)=1.
		\end{array}\right.
	\end{equation}

		\begin{equation}\label{u_xi}
			u_\xi=\frac{q}{2}\sin v.
		\end{equation}

		\begin{equation}\label{bdd_u}
			\sup_{\xi\in\mathbb{R}}|u^{2}(t,\xi)| \leq 2\int_\mathbb{R}|uu_\xi|d\xi  \leq2\int_\mathbb{R}|u||\sin\frac{v}{2}\cos\frac{v}{2}|qd\xi\leq \frac{1}{\sqrt{\mu}}E_0.
		\end{equation}

		\begin{equation}\label{partial_t y(t,xi)}
			y_t(t,\xi)=\frac{s}{\mu}u(t,\xi)-\frac{\eta}{\mu},\quad y(0,\xi)=y_0(\xi).
		\end{equation}

		\begin{equation}\label{y_xi}
			y_\xi=q\cos^2\frac{v}{2}.
		\end{equation}

		\begin{equation}\label{bdd_Px_Linfty}
			\|P_x\|_{L^\infty}\leq C\left(E_0^{\frac{1}{2}}+E_0+E_0^{\frac{m+1}{2}}\right).
		\end{equation}

		\begin{equation}\label{bdd_y(t,xi)}
			y_0(\xi)-C\left(1+E_0^{\frac{1}{2}}\right)t\leq y(t,\xi)\leq y_0(\xi)+C\left(1+E_0^{\frac{1}{2}}\right)t,
		\end{equation}

	\begin{lemm}\label{lemma_N}
		Let $s\neq0$, 
		$$\mathcal{N}\triangleq\left\{t\geq0:\mathrm{meas}\{\xi\in \mathbb{R}:v(t,\xi)=-\pi\}>0\right\}.$$
		Then we have
		$$\mathrm{meas}(\mathcal{N})=0.$$
	\end{lemm}

	\begin{theo}
		Let $s\neq0,\ u_0\in H^1(\mathbb{R})$. Then the Cauchy problem \eqref{eq1} has a global conservative solution in the sense of Definition \ref{conservative}.
	\end{theo}

	\begin{proof}
		Given a global solution $(u,v,q)$ of \eqref{semi-linear system}-\eqref{semi-linear system_initial data}, we set
		\begin{equation}\label{def_u(t,x)}
			u(t,x)\triangleq u(t,\xi),\quad\mathrm{if}\quad x=y(t,\xi),
		\end{equation}
		with $y(t,\xi)$ given by \eqref{partial_t y(t,xi)}. Now we check that \eqref{def_u(t,x)} is well defined. From \eqref{y0} and \eqref{bdd_y(t,xi)}, we see
		$$\lim_{\xi\to\pm\infty}y(t,\xi)=\pm\infty.$$
		Next, \eqref{y_xi} implies that the map $\xi\mapsto y(t,\xi)$ is non-decreasing. If $\xi_1<\xi_2$ but $y(t,\xi_1)=y(t,\xi_2)$, then
		$$\int_{\xi_1}^{\xi_2}q(t,s)\cos^2\frac{v(t,s)}{2}ds=\int_{\xi_1}^{\xi_2}y_\xi(t,s)ds=0,$$
		which yields $\cos\frac{v}{2}=0$ throughout the interval of integration. Thus, \eqref{u_xi} gives
		$$u(t,\xi_2)-u(t,\xi_1)=\int_{\xi_1}^{\xi_2}\frac{q(t,s)}{2}\sin v(t,s)ds=0,$$
		which proves that \eqref{def_u(t,x)} is well defined for all $t\geq0$ and $x\in\mathbb{R}$.
		
		According to the definition \eqref{def_u(t,x)}, we have
		\begin{equation*}
			u_x(t,x)=\frac{\sin v(t,\xi)}{1+\cos v(t,\xi)}\quad\mathrm{if}\quad x=y(t,\xi),\quad\cos v(t,\xi)\neq-1,
		\end{equation*}
		which, together with \eqref{y_xi}, implies
		\begin{align*}
			&\int_{\mathbb{R}}u^{2}(t,x)+\mu u_{x}^{2}(t,x)dx \\
			=&\int_{\{\cos v>-1\}}\left(u^2(t,\xi)\cos^2\frac{v(t,\xi)}{2}+\mu\sin^2\frac{v(t,\xi)}{2}\right)q(t,\xi)d\xi \\
			\leq& E_{0}.
		\end{align*}
		Since $H^1\hookrightarrow C^{0,\frac{1}{2}}$, we see that $u$ is uniformly H\"{o}lder continuous with exponent $\frac{1}{2}$ as a function of $x$. According to \eqref{semi-linear system} and \eqref{bdd_Px_Linfty}, the map $t\mapsto u(t,y(t))$ is uniformly Lipschitz continuous along each characteristic curve $t\mapsto y(t)$. Hence, $u=u(t,x)$ is globally H\"{o}lder continuous.
		
		Now we claim that the map $t\mapsto u(t,\cdot)$ is Lipschitz continuous with values in $L^2(\mathbb{R})$. Indeed, for any interval $[\tau,\tau+h]$ and a given point $x$, we choose $\xi\in\mathbb{R}$ such that the characteristic $t\mapsto y(t,\xi)$ passes through $(\tau,x)$. It follows from \eqref{semi-linear system} and \eqref{bdd_u} that 
		$$\|y_t\|_{L^\infty}\leq\frac{|s|E_0^{\frac{1}{2}}}{\mu^{\frac{5}{4}}}+\frac{|\eta|}{\mu}\triangleq\widetilde{C}$$
		and 
		\begin{align*}
			&\left|u(\tau+h,x)-u(\tau,x)\right|\\
			\leq&\left|u(\tau+h,x)-u\left(\tau+h,y(\tau+h,\xi)\right)\right|  +\left|u\left(\tau+h,y(\tau+h,\xi)\right)-u(\tau,x)\right| \\
			\leq&\sup_{|y-x|\leq \widetilde{C}h}\left|u(\tau+h,y)-u(\tau+h,x)\right|+\int_{\tau}^{\tau+h}\left|P_{x}(t,\xi)\right|dt.
		\end{align*}
		Integrating over $\mathbb{R}$ yields
		\begin{align*}
			\int_\mathbb{R}\left|u(\tau+h,x)-u(\tau,x)\right|^2dx\leq&2\int_\mathbb{R}\left(\int_{x-\widetilde{C}h}^{x+\widetilde{C}h}\left|u_x(\tau+h,y)\right|dy\right)^2dx\\
			&+2\int_\mathbb{R}\left(\int_\tau^{\tau+h}\left|P_x(t,\xi)\right|dt\right)^2q(\tau,\xi)\cos^2\frac{v(\tau,\xi)}{2}d\xi\\
			\leq&8\widetilde{C}^2h^2\|u_x(\tau+h)\|_{L^2}^2+2h\|q(\tau)\|_{L^\infty}\int_\tau^{\tau+h}\|P_x(t)\|_{L^2}^2dt\\
			\leq&Ch^2,
		\end{align*}
		thus proving the claim.
		
		Since $L^2(\mathbb{R})$ is a reflexive space, the left-hand side of \eqref{u} is well defined for a.e. $t\geq0$ and the right-hand side of \eqref{u} also lies in $L^2(\mathbb{R})$ for a.e. $t\geq0$. From \eqref{semi-linear system}, we observe that for any $t>0$ and $\xi\in\mathbb{R}$,
		$$\frac{d}{dt}u\left(t,y(t,\xi)\right)=-P_x(t,\xi),$$
		where $P_x$ is the function defined at \eqref{Px}. Moreover, for every $t\notin\mathcal{N}$, the map $\xi\mapsto y(t,\xi)$ is one-to-one and a
		change of variables yields
		\begin{align*}
			P_x(t,\xi) 
			=&\frac{1}{2\mu}\left(\int_{y(t,\xi)}^{\infty}-\int_{-\infty}^{y(t,\xi)}\right)\exp\left\{-\frac{1}{\sqrt{\mu}}\left|y(t,\xi)-x\right|\right\}\\
			&\cdot\left(-\left(\frac{A}{2}+\frac{s}{2\mu}\right)u^2+\frac{s}{2}u_x^2-\frac{B}{m+1}u^{m+1}+\left(2k+\frac{\eta}{\mu}\right)u\right)(t,x)dx\\
			=&P_x(t,y(t,\xi)).
		\end{align*}
		Since Lemma \ref{lemma_N} guarantees $\mathrm{meas}(\mathcal{N})=0$, the identity \eqref{u} holds for almost every $t\geq0$.
		
		Finally, we consider the balance law. Let $\left\{\mu_{(t)},\ t\geq0\right\}$ be the bounded Radon measures defined by
		\begin{equation}
			\mu_{(t)}\left([a,b]\right)\triangleq\int_{\{\xi;y(t,\xi)\in[a,b]\}}\sin^2\frac{v(t,\xi)}{2}\cdot q(t,\xi)d\xi.
		\end{equation}
		If $t\notin\mathcal{N}$, then the measure $\mu_{(t)}$ is absolutely continuous, possessing density $u_x^2$ with respect to the Lebesgue measure. Moreover, for $t\in\mathcal{N}$, $\mu_{(t)}$ is the weak limit of $\mu_{(t^\prime)}$, as $t^\prime\to t,\ t^\prime\notin\mathcal{N}$. It follows from \eqref{semi-linear system} and \eqref{u_xi} that for any test function $\varphi\in\mathcal{C}_c^1(\mathbb{R}^+\times\mathbb{R})$,
		\begin{align*}
			&-\int_{0}^{\infty}\int_\mathbb{R}\varphi_t+\varphi_x\left(\frac{s}{\mu}u-\frac{\eta}{\mu}\right)d\mu_{(t)}dt\\
			=&-\int_{0}^{\infty}\int_\mathbb{R}\partial_t(\varphi(t,y(t,\xi)))\cdot\sin^2\frac{v(t,\xi)}{2}\cdot q(t,\xi)d\xi dt\\
			=&\frac{1}{\mu}\int_{0}^{\infty}\int_\mathbb{R}\varphi(t,y(t,\xi))\cdot\left[\left(-\left(\frac{A}{2}+\frac{s}{2\mu}\right)u^2-\frac{B}{m+1}u^{m+1}+\left(2k+\frac{\eta}{\mu}\right)u-P\right)\sin v\cdot q\right](t,\xi) d\xi dt\\
			&+\int_{\mathbb{R}}\varphi(0,y(0,\xi))\cdot\sin^2\frac{v(0,\xi)}{2}\cdot q(0,\xi)d\xi\\
			=&\frac{2}{\mu}\int_{0}^{\infty}\int_\mathbb{R}\left(-\left(\frac{A}{2}+\frac{s}{2\mu}\right)u^2-\frac{B}{m+1}u^{m+1}+\left(2k+\frac{\eta}{\mu}\right)u-P\right)u_x\varphi  dxdt+\int_{\mathbb{R}}\varphi(0,x)u_{0x}^{2}dx.
		\end{align*}	
		%Since Lemma \ref{lemma_N} gives $\mathrm{meas}(\mathcal{N})=0$, \eqref{ux2} yields the balance law \eqref{eqw}, which can be formulated more precisely as
		
		Hence, we complete the proof that $u$ is a global conservative solution of \eqref{eq1} in the sense of Definition \ref{conservative}.
	\end{proof}

\end{document}
